\chapter{Networking Devices}\label{networking-devices}

\section{What is a network?}\label{what-is-a-network}

A computer network is a digital telecommunications network allows NODES
to share RESOURCES.

A CLIENT is a device that accesses a service made available by a SERVER.

A SERVER is a device that provides functions or services for CLIENTS.

\begin{itemize}
\tightlist
\item
  Note : The same device can be a CLIENT in some situations and a SERVER
  in other situations. Ex: A Peer-to-Peer network.
\end{itemize}

SWITCHES (Level 2):

\begin{itemize}
\tightlist
\item
  provide connectivity to hosts within the same LAN (Local Area Network)
\item
  Have many network interfaces/ports for End Hosts to connect to.
\item
  DO NOT provide connectivity between LANs/over the Internet.
\end{itemize}

ROUTERS (Level 3):

\begin{itemize}
\tightlist
\item
  have fewer network interfaces than switches.
\item
  are used to provide connectivity BETWEEN LANs.
\item
  are used to send data over the Internet.
\end{itemize}

FIREWALL (Can be Level 3,4, and 7):

\begin{itemize}
\tightlist
\item
  Firewalls are specialty hardware network security devices that control
  network traffic entering/exiting your network.
\item
  Can be places ``inside'' or ``outside'' the network.
\item
  Monitor and control network traffic based on configured rules.
\item
  Are known as ``Next-Generation Firewalls'' when they include more
  modern and advanced filtering capabilities.
\item
  Host-based firewalls are software applications that filter traffic
  entering and exiting a host machine, like a PC.
\end{itemize}
